\scanpage{40}
\mspage{127}

1. \uline{Hot-fibrations} \uline{dans} (Cat)

\deleted{\unsure{Proposition 1}} \deleted{\illegible{1}} \inserted{\uline{Théorème 1} :} Soit $f:X\longrightarrow Y$ dans Cat. Conditions équivalentes :

a) Pour tout $Y''\xrightarrow{p}Y'$ sur $Y$, d'où $X''=X\times_Y Y''\xrightarrow{q}X'=X\times_Y Y'$,
si $p$ est une Hot-équivalence, $q$ l'est aussi.

a$'$) Comme a, avec $Y''=e$, $Y'$ ayant un objet final

\deleted{a$''$) Comme \unsure{a)}, avec $Y''=e$, $Y'$ ayant un objet initial}

% slanted
\margin{garder a$''$) et b$''$), conditions duales de a$'$), b$'$)}

b$'$) Pour \unsure{tout} $y\in\mathrm{Ob}\,Y$, $X_y\to X_{/y}$ est une Hot équivalence,
pour \unsure{toute} $y\to y'$ dans $Y$, $X_{/y}\to X_{/y'}$ est une Hot-équivalence,
enfin ces \deleted{\unsure{deux assertions}} conditions sont aussi vérifiées
pour $f':X'\to Y'$, pour tout chgt. de base $Y'\unsure{\text{---}|\to}Y$
dans Cat. (En fait, il suffit de le demander pour la première des
deux conditions, \unsure{nommément|concernant} $X_y\to X_{/y}$.)

\deleted{b$''$) Dual de b$'$, avec $X_y\to {}_{y\backslash}X$ et $ {}_{y'\backslash}X\to {}_{y\backslash}X$}

c) Pour tout faisceau d'ens loc. constant \inserted{$F$} sur $X$,
$f_*(F)$ est loc. constant et se calcule ``fibre par fibre'',
de $=$ \inserted{pour} $R^1f_*(F)$ si $F$ est un faisceau en \unsure{groupes|groupe},
et pour les $\mathrm R^i f_*(F)$ si $F$ est un faisceau abélien. De
plus, on veut que ces conditions restent satisfaites
après tout chgt. de base.

% slanted
\margin{\textcircled{c$'$} version homologique $1^\circ$) en \deleted{\unsure{termes}} d'images \unsure{directes} de groupoïdes $\unsure{f_*|f_!}(F)$, loc. constants, et $2^\circ$) en termes des $\deleted{\unsure{H_i}}\unsure{R^i|R^1} f_!(L)$, pour $L$ faisceau abélien loc. constant.}

\uline{Dém} On a les implications $\pm$ tautologiques
\diagram{\begin{tikzcd}[ampersand replacement=\&,row sep=2pt,column sep=small,execute at end picture={\draw[decorate,decoration={zigzag,segment length=3pt,amplitude=3pt}] (aa.south east) -- (aa2.west) -- (bb2.west);}]
\& a'\arrow[r,Rightarrow]\& b'\arrow[dr,Rightarrow]\&\\
|[alias=aa]| a\arrow[ur,Rightarrow]\arrow[dr,Rightarrow]\&\&\&c,\\
\&|[alias=aa2]| \deleted{\unsure{a''}}\arrow[r,Rightarrow]\&|[alias=bb2]| \deleted{b''}\&
\end{tikzcd}}
et il reste à prouver que $c\Longrightarrow a$.
% slanted
\margin{Comme a est autoduale, étant équivalente à a$'$, b$'$, elle est aussi équivalente aux \uline{conditions duales} de celles-ci a$''$ et b$''$. (cf cor. 2, page 75)}
On est ramené à prouver que si c) est satisfait, et si
\diagram{\begin{tikzcd}[ampersand replacement=\&,row sep=small,column sep=small]
X\arrow[d,"f"']\& X'\arrow[l,"q"']\arrow[d,"f'"]\\
Y\& Y'\arrow[l,"p"]
\end{tikzcd}}
cartésien dans Cat, si $p$ est une Hot-équiv., $q$ aussi.

On utilise le critère d'Artin-Mazur.

$1^\circ$) Si $F$ est un faisceau d'ens. \inserted{loc.} constant sur $Y$,
il faut prouver que $H^0(X,F)\simeq H^0(X',q^*F)$.

% archive sheet number: 40
